\par\Needspace{5\baselineskip}
\begin{center}\SourceLeaf{418}3.\end{center}
\noindent\hspace*{1.2em}Suppose we refer the position of points in the plane to two perpendicular coordinate axes to which the directions $0$ and $90^{0}$ correspond, so that the first coordinate represents the distance from the second axis measured in the direction of the first, while the second coordinate represents the distance from the first axis measured in the direction of the second. If further $x,y$ denote the coordinates of an arbitrary point on the curved line, $s$ its length counted from an arbitrary initial point up to that point, $\varphi$ the direction of the tangent there, and $r$ the radius of curvature, then
\[
\begin{aligned}
dx&=\cos\varphi.ds\\
dy&=\sin\varphi.ds\\
r&=\frac{ds}{d\varphi}.
\end{aligned}
\]

If the form of the curved line is represented by the equation $V=0$, where $V$ is a function of $x$ and $y$, and we put
\[
dV=pdx+qdy,
\]
then along the curved line we have
\[
pdx+qdy=0
\]
and hence
\[
p\cos\varphi+q\sin\varphi=0,
\]
from which
\[
\operatorname{tang}\varphi=-\frac pq.
\]
Further,
\[
\cos\varphi.dp+\sin\varphi.dq-(p\sin\varphi-q\cos\varphi)d\varphi=0.
\]
If, then, by a well-known theorem, we put
\[
\begin{aligned}
dp&=Pdx+Qdy\\
dq&=Qdx+Rdy,
\end{aligned}
\]
